% Part: methods
% Chapter: induction
% Section: inductive-definitions

\documentclass[../../../include/open-logic-section]{subfiles}

\begin{document}

\olfileid[hi]{mth}{ind}{idf}

\olsection{आगमनात्मक परिभाषाएँ}

तर्कशास्त्र में हम वस्तुओं के प्रकार प्रायः \emph{आगमनात्मक रूप से}
परिभाषित करते हैं, अर्थात् यह बताने वाले नियम देकर कि परिभाषित किए जाने
वाले प्रकार की वस्तु किसे माना जाए और उस प्रकार की पुरानी वस्तुओं से नई
वस्तुएँ कैसे मिलें। उदाहरण के लिए हम किसी भाषा के पदों और !!{formula}ों
जैसी चिह्नों की विशिष्ट प्रकार की शृंखलाएँ प्रायः आगमन से परिभाषित करते हैं।
सरल उदाहरण के लिए $\mathrm{a}$, $\mathrm{b}$, $\mathrm{c}$,
$\mathrm{d}$ अक्षरों, चिह्न $\circ$ और कोष्ठकों $[$ तथा $]$ से बनी
शृंखलाओं पर विचार करें, जैसे ``$[[\mathrm{c}\circ\mathrm{d}][$'',
``$[\mathrm{a}[]\circ]$'', ``$\mathrm{a}$'' अथवा
``$[[\mathrm{a}\circ\mathrm{b}]\circ\mathrm{d}]$''। आपको शायद लगे कि
पहली दो शृंखलाओं में कुछ ``गलत'' है: पहली में कोष्ठक बिल्कुल ``संतुलित''
नहीं हैं, और आपको लग सकता है कि ``$\circ$'' को ऐसी अभिव्यक्तियाँ ``जोड़नी''
चाहिए जो स्वयं सार्थक हों। तीसरी और चौथी शृंखलाएँ बेहतर दिखती हैं: प्रत्येक
``$[$'' के लिए समापन ``$]$'' है (यदि कोई कोष्ठक हैं), और प्रत्येक $\circ$
के दोनों ओर कोष्ठकों के युग्म से घिरी ``सुंदर'' अभिव्यक्तियाँ मिलती हैं।

हम सटीक रूप से बताना चाहते हैं कि ``सुंदर पद'' किसे माना जाए। सबसे पहले
प्रत्येक अक्षर अपने-आप में सुंदर है। जो केवल एक अक्षर नहीं है, उसे
``$[t\circ s]$'' रूप का होना चाहिए, जहाँ $s$ और $t$ स्वयं सुंदर हों। इसके
विपरीत यदि $t$ और $s$ सुंदर हैं, तो उनके बीच $\circ$ रखकर और उन्हें
कोष्ठकों के युग्म से घेरकर नया सुंदर पद बना सकते हैं। इन संक्रियाओं का उपयोग
सुंदर पदों के समुच्चय को \emph{परिभाषित} करने में कर सकते हैं। यह एक
\emph{आगमनात्मक परिभाषा} है।

\begin{defn}[सुंदर पद]
  \emph{सुंदर पदों} का समुच्चय इस प्रकार आगमनात्मक रूप से परिभाषित है:
  \begin{enumerate}
  \item कोई भी अक्षर $\mathrm{a}$, $\mathrm{b}$, $\mathrm{c}$,
    $\mathrm{d}$ एक सुंदर पद है।
  \item यदि $s_1$ और $s_2$ सुंदर पद हैं, तो $[s_1\circ s_2]$ भी सुंदर पद
    है।
  \item और कुछ भी सुंदर पद नहीं है।
  \end{enumerate}
\end{defn}

यह परिभाषा बताती है कि कोई वस्तु सुंदर पद ठीक तभी मानी जाती है जब उसे (1)
और (2) की शर्तों के अनुसार परिमित चरणों में निर्मित किया जा सके। पहले चरण
में हम केवल अलग-अलग अक्षरों से बने सभी सुंदर पद बनाते हैं, अर्थात्,
\[
\mathrm{a}, \mathrm{b}, \mathrm{c}, \mathrm{d}
\]
दूसरे चरण में हम निर्मित पदों पर (2) लागू करते हैं। हमें मिलेंगे
\[
[\mathrm{a} \circ \mathrm{a}], [\mathrm{a} \circ \mathrm{b}],
[\mathrm{b} \circ \mathrm{a}], \dots, [\mathrm{d} \circ \mathrm{d}]
\]
दो अक्षरों के सभी संयोजनों के लिए। तीसरे चरण में अब तक बनाए किन्हीं दो सुंदर
पदों पर (2) फिर लागू करते हैं। हमें
$[\mathrm{a}\circ[\mathrm{a}\circ\mathrm{a}]]$ जैसे नए सुंदर पद मिलते हैं
---जहाँ $t$ चरण~1 का $\mathrm{a}$ और $s$ चरण~2 का
$[\mathrm{a}\circ\mathrm{a}]$ है---और चरण~2 के दो पदों
$[\mathrm{b}\circ\mathrm{c}]$ तथा $[\mathrm{d}\circ\mathrm{b}]$ से बना
$[[\mathrm{b}\circ\mathrm{c}]\circ[\mathrm{d}\circ\mathrm{b}]]$। और इसी
प्रकार आगे। खंड (3) ऐसी किसी वस्तु को सुंदर पदों के समुच्चय में घुसने से
रोकता है जो इस ढंग से निर्मित नहीं हुई।

ध्यान दें कि हमने अभी यह सिद्ध नहीं किया कि चिह्नों की हर शृंखला जो सुंदर
``लगती'' है, इस परिभाषा के अनुसार सुंदर है। किंतु यह स्पष्ट होना चाहिए कि
जो कुछ हम निर्मित कर सकते हैं वह वास्तव में सुंदर ``लगता'' है: कोष्ठक
संतुलित हैं और $\circ$ ऐसे भागों को जोड़ता है जो स्वयं सुंदर हैं।

आगमनात्मक परिभाषाओं की प्रमुख विशेषता यह है कि यदि आप सभी सुंदर पदों के
बारे में कुछ सिद्ध करना चाहते हैं, तो परिभाषा बताती है कि किन स्थितियों पर
विचार करना होगा। उदाहरणतः यदि बताया जाए कि $t$ एक सुंदर पद है, तो आगमनात्मक
परिभाषा बताती है कि $t$ कैसा दिख सकता है: $t$ कोई अक्षर हो सकता है, या सुंदर
पदों के किसी युग्म $s_1$ और $s_2$ के लिए $[s_1\circ s_2]$ हो सकता है। खंड
(3) के कारण केवल यही संभावनाएँ हैं।

आगमनात्मक रूप से परिभाषित समुच्चय की सभी वस्तुओं के बारे में दावे सिद्ध
करते समय आगमन का प्रबल रूप विशेष रूप से महत्त्वपूर्ण हो जाता है। उदाहरणतः
मान लें हम सिद्ध करना चाहते हैं कि लंबाई $n$ के प्रत्येक सुंदर पद में $[$
की संख्या $<n/2$ है। इसे सभी $n$ के बारे में दावा माना जा सकता है: प्रत्येक
$n$ के लिए लंबाई $n$ के किसी भी सुंदर पद में $[$ की संख्या $<n/2$ है।

\begin{prop}
  किसी भी $n$ के लिए लंबाई $n$ के सुंदर पद में $[$ की संख्या $<n/2$ है।
\end{prop}

\begin{proof}
इस परिणाम को (प्रबल) आगमन से सिद्ध करने के लिए हमें दिखाना होगा कि निम्न
सशर्त दावा सत्य है:
\begin{quote}
  यदि प्रत्येक $l<k$ के लिए लंबाई $l$ के किसी भी सुंदर पद में $[$ की
  संख्या $<l/2$ है, तो लंबाई $k$ के किसी भी सुंदर पद में $[$ की संख्या
  $<k/2$ है।
\end{quote}
इस सशर्त को दिखाने के लिए उसके पूर्ववर्ती को सत्य मानें, अर्थात् मानें कि
किसी भी $l<k$ के लिए लंबाई $l$ के सुंदर पदों में $[$ की संख्या $<l/2$ है।
इस मान्यता को हम आगमन परिकल्पना कहते हैं। हम दिखाना चाहते हैं कि लंबाई $k$
के सुंदर पदों के लिए भी यही सत्य है।

अतः मान लें $t$ लंबाई $k$ का सुंदर पद है। चूँकि सुंदर पद आगमनात्मक रूप से
परिभाषित हैं, दो स्थितियाँ हैं: (1) $t$ स्वयं एक अक्षर है, अथवा (2) किन्हीं
सुंदर पदों $s_1$ और $s_2$ के लिए $t=[s_1\circ s_2]$।
\begin{enumerate}
\item $t$ एक अक्षर है। तब $k=1$, और $t$ में $[$ की संख्या $0$ है। चूँकि
$0<1/2$, दावा सत्य है।
\item किन्हीं सुंदर पदों $s_1$ और $s_2$ के लिए $t=[s_1\circ s_2]$। $l_1$
  को $s_1$ की लंबाई और $l_2$ को $s_2$ की लंबाई मानें। तब $t$ की लंबाई
  $k=l_1+l_2+3$ है ($s_1$ और $s_2$ की लंबाइयाँ और तीन चिह्न $[$,
  $\circ$, $]$)। चूँकि $l_1+l_2+3$ सदैव $l_1$ से बड़ा है, $l_1<k$। इसी
  प्रकार $l_2<k$। इसका अर्थ है कि आगमन परिकल्पना $s_1$ और $s_2$ पर लागू
  होती है: $s_1$ में $[$ की संख्या $m_1<l_1/2$, और $s_2$ में $[$ की
  संख्या $m_2<l_2/2$।

  $t$ में $[$ की संख्या $s_1$ में $[$ की संख्या, और $s_2$ में $[$ की
  संख्या, और $1$ का योग है, अर्थात् $m_1+m_2+1$। चूँकि $m_1<l_1/2$ और
  $m_2<l_2/2$, हमारे पास है:
  \[
  m_1 + m_2 + 1 < \frac{l_1}{2} + \frac{l_2}{2} + 1 = \frac{l_1+l_2+2}{2} < \frac{l_1+l_2+3}{2} = k/2.
  \]
\end{enumerate}
प्रत्येक स्थिति में हमने (आगमन परिकल्पना के आधार पर) दिखाया कि $t$ में $[$
की संख्या $<k/2$ है। प्रबल आगमन से प्रतिज्ञप्ति निष्पन्न होती है।
\end{proof}

\begin{prob}
  अतिसुंदर पदों का समुच्चय इस प्रकार परिभाषित करें
  \begin{enumerate}
  \item कोई भी अक्षर $\mathrm{a}$, $\mathrm{b}$, $\mathrm{c}$,
    $\mathrm{d}$ अतिसुंदर पद है।
  \item यदि $s$ अतिसुंदर पद है, तो $[s]$ भी अतिसुंदर पद है।
  \item यदि $s_1$ और $s_2$ अतिसुंदर पद हैं, तो $[s_1\circ s_2]$ भी
    अतिसुंदर पद है।
  \item और कुछ भी अतिसुंदर पद नहीं है।
  \end{enumerate}
  दिखाइए कि लंबाई $n$ के अतिसुंदर पद $t$ में $[$ की संख्या
  $\le n/2 +1$ है।
\end{prob}

\end{document}
